\documentclass[12pt,a4paper]{article}

\input{/home/tabaka/Documents/GitHub/Defs/defs.tex}

\setcounter{zadaniaRozdzial}{3}
\setcounter{section}{0}

\begin{document}
	
	\begin{center}
		\Large Wymagania edukacyjne z matematyki w klasie 2 Liceum Ogólnokształcącego na poszczególne oceny.
	\end{center}
	
	Wymagania edukacyjne na poszczególne oceny są stopniowane.
	Uzyskanie danej oceny wymaga spełnienia wymagań przewidzianych
	dla wszystkich ocen niższych, z uwzględnieniem wymagań
	określonych dla danej oceny.
	
	\section{Funkcja kwadratowa}
	
	\begin{longtable}{|p{16cm}|}
		\hline
		
		\textbf{Wymagania edukacyjne na ocenę dopuszczającą} \\ \hline
		
		\textbullet\ szkicuje wykres funkcji kwadratowej $f(x)=ax^2$, gdzie $a\neq 0$, i odczytuje z wykresu jej własności; \\
		
		\textbullet\ szkicuje wykres funkcji kwadratowej $f(x)=ax^2+q$, gdzie $a\neq 0$, i odczytuje z wykresu jej własności; \\
		
		\textbullet\ szkicuje wykres funkcji kwadratowej $f(x)=a(x-p)^2$, gdzie $a\neq 0$, i odczytuje z wykresu jej własności; \\
		
		\textbullet\ podaje wzór funkcji kwadratowej w postaci ogólnej i kanonicznej; \\
		
		\textbullet\ oblicza współrzędne wierzchołka paraboli i wyznacza równanie jej osi symetrii; \\
		
		\textbullet\ przekształca postać kanoniczną funkcji kwadratowej do postaci ogólnej; \\
		
		\textbullet\ rozwiązuje niepełne równanie kwadratowe metodą wyłączania wspólnego czynnika przed nawias lub z zastosowaniem wzorów skróconego mnożenia; \\
		
		\textbullet\ określa liczbę pierwiastków równania kwadratowego w zależności od znaku wyróżnika; \\
		
		\textbullet\ rozwiązuje równanie kwadratowe w prostych przypadkach, stosując wzory na pierwiastki; \\
		
		\textbullet\ odczytuje miejsca zerowe funkcji kwadratowej zapisanej w postaci iloczynowej; \\
		
		\textbullet\ rozwiązuje nierówność kwadratową w prostych przypadkach. \\ \hline
		
		\textbf{Wymagania edukacyjne na ocenę dostateczną} \\ \hline
		
		\textbullet\ szkicuje wykres funkcji kwadratowej $f(x)=a(x-p)^2+q$, gdzie $a\neq 0$, i odczytuje z wykresu jej własności; \\
		
		\textbullet\ przekształca postać ogólną funkcji kwadratowej do postaci kanonicznej, korzystając ze wzorów na współrzędne wierzchołka paraboli, a następnie szkicuje wykres funkcji i opisuje jej własności; \\
		
		\textbullet\ wyznacza wzór funkcji kwadratowej w postaci ogólnej, gdy dane są współrzędne wierzchołka paraboli i innego punktu jej wykresu; \\
		
		\textbullet\ rozwiązuje równanie kwadratowe, stosując wzory na pierwiastki; \\
		
		\textbullet\ interpretuje geometrycznie rozwiązania równania kwadratowego w zależności od współczynnika $a$ i wyróżnika $\Delta$; \\
		
		\textbullet\ wyznacza algebraicznie współrzędne punktów przecięcia paraboli z osiami układu współrzędnych; \\
		
		\textbullet\ przedstawia trójmian kwadratowy w postaci iloczynowej, jeśli taka postać istnieje w rozważanym zbiorze liczb; \\
		
		\textbullet\ rozwiązuje nierówności kwadratowe; \\
		
		\textbullet\ stosuje pojęcia najmniejszej i największej wartości funkcji oraz wyznacza w prostych przypadkach najmniejszą i największą wartość funkcji kwadratowej w przedziale domkniętym; \\
		
		\textbullet\ analizuje zadanie tekstowe, zapisuje odpowiednie równanie, nierówność lub wzór funkcji kwadratowej opisujące daną zależność i znajduje w prostych przypadkach rozwiązanie spełniające warunki zadania. \\ \hline
		
		\textbf{Wymagania edukacyjne na ocenę dobrą} \\ \hline
		
		\textbullet\ rozwiązuje równania i nierówności kwadratowe w trudniejszych przypadkach; \\
		
		\textbullet\ wykorzystuje postać iloczynową funkcji kwadratowej do rozwiązywania zadań; \\
		
		\textbullet\ stosuje nierówności kwadratowe do wyznaczania dziedziny funkcji zapisanej za pomocą pierwiastka; \\
		
		\textbullet\ wyznacza w trudniejszych przypadkach najmniejszą i największą wartość funkcji kwadratowej w przedziale domkniętym, korzystając z jej własności; \\
		
		\textbullet\ stosuje własności funkcji kwadratowej do rozwiązywania zadań optymalizacyjnych. \\ \hline
		
		\textbf{Wymagania edukacyjne na ocenę bardzo dobrą} \\ \hline
		
		\textbullet\ zaznacza na osi liczbowej sumę, iloczyn i różnicę zbiorów rozwiązań dwóch nierówności kwadratowych; \\
		
		\textbullet\ wykorzystuje postać iloczynową funkcji kwadratowej do rozwiązywania trudniejszych zadań; \\
		
		\textbullet\ stosuje własności funkcji kwadratowej do rozwiązywania trudniejszych zadań optymalizacyjnych; \\
		
		\textbullet\ rozwiązuje zadania tekstowe o podwyższonym stopniu trudności, stosując równania kwadratowe. \\ \hline
		
		\textbf{Wymagania edukacyjne na ocenę celującą} \\ \hline
		
		\textbullet\ wyprowadza wzory na pierwiastki trójmianu kwadratowego; \\
		
		\textbullet\ udowadnia zależności między współczynnikami funkcji kwadratowej w zadaniach o podwyższonym stopniu trudności; \\
		
		\textbullet\ rozwiązuje zadania o znacznym stopniu trudności dotyczące funkcji kwadratowej. \\ \hline
		
	\end{longtable}
	
	\section{Wielomiany}
	
	\begin{longtable}{|p{16cm}|}
		\hline
		
		\textbf{Wymagania edukacyjne na ocenę dopuszczającą} \\ \hline
		
		\textbullet\ podaje przykład wielomianu, określa jego stopień i podaje wartości jego współczynników; \\
		
		\textbullet\ zapisuje wielomian określonego stopnia o danych współczynnikach; \\
		
		\textbullet\ zapisuje wielomian w postaci uporządkowanej; \\
		
		\textbullet\ oblicza wartość wielomianu dla danego argumentu wymiernego; \\
		
		\textbullet\ wyznacza sumę, różnicę i iloczyn wielomianów oraz określa ich stopnie; \\
		
		\textbullet\ określa stopień iloczynu wielomianów bez wykonywania mnożenia; \\
		
		\textbullet\ przekształca wyrażenia algebraiczne, stosując wzory skróconego mnożenia dotyczące wielomianów stopnia drugiego; \\
		
		\textbullet\ rozwiązuje równania wielomianowe, w których wielomian jest zapisany w postaci iloczynu czynników liniowych. \\ \hline
		
		\textbf{Wymagania edukacyjne na ocenę dostateczną} \\ \hline
		
		\textbullet\ oblicza wartość wielomianu dla danego argumentu niewymiernego; \\
		
		\textbullet\ sprawdza, czy dany punkt należy do wykresu wielomianu; \\
		
		\textbullet\ wyznacza współczynnik przy najwyższej potędze oraz wyraz wolny iloczynu wielomianów bez wykonywania mnożenia; \\
		
		\textbullet\ przekształca wyrażenia algebraiczne, stosując wzory skróconego mnożenia dotyczące wielomianów stopnia drugiego; \\
		
		\textbullet\ rozwiązuje równania wielomianowe, w których wielomian jest zapisany w postaci iloczynu czynników stopnia co najwyżej drugiego; \\
		
		\textbullet\ opisuje wielomianem zależności przedstawione w zadaniu i wyznacza dziedzinę tego wielomianu w prostych przypadkach. \\ \hline
		
		\textbf{Wymagania edukacyjne na ocenę dobrą} \\ \hline
		
		\textbullet\ wyznacza współczynniki wielomianu spełniającego podane warunki; \\
		
		\textbullet\ oblicza wartość wielomianu dwóch lub trzech zmiennych dla podanych argumentów; \\
		
		\textbullet\ wyznacza iloczyn wielomianów wielu zmiennych; \\
		
		\textbullet\ rozwiązuje równania wielomianowe w trudniejszych przypadkach, gdy wielomian jest zapisany w postaci iloczynu czynników stopnia co najwyżej drugiego; \\
		
		\textbullet\ sprawdza, czy rozwiązania równania wielomianowego spełniają podane warunki; \\
		
		\textbullet\ opisuje wielomianem pole powierzchni prostopadłościanu i określa dziedzinę tego wielomianu. \\ \hline
		
		\textbf{Wymagania edukacyjne na ocenę bardzo dobrą} \\ \hline
		
		\textbullet\ określa stopień sumy wielomianów w zależności od wartości parametru; \\
		
		\textbullet\ rozwiązuje zadania tekstowe, wykorzystując działania na wielomianach i równania wielomianowe. \\ \hline
		
		\textbf{Wymagania edukacyjne na ocenę celującą} \\ \hline
		
		\textbullet\ rozwiązuje zadania o znacznym stopniu trudności dotyczące wielomianów. \\ \hline
		
	\end{longtable}
	
	\section{Funkcje wymierne}
	
	\begin{longtable}{|p{16cm}|}
		\hline
		
		\textbf{Wymagania edukacyjne na ocenę dopuszczającą} \\ \hline
		
		\textbullet\ szkicuje w prostych przypadkach wykres funkcji $f(x)=\frac{a}{x}$, gdzie $a\neq 0$, oraz podaje jej własności: dziedzinę, zbiór wartości i przedziały monotoniczności; \\
		
		\textbullet\ przesuwa wykres funkcji $f(x)=\frac{a}{x}$, gdzie $a\neq 0$, wzdłuż osi $OX$ lub osi $OY$, podaje własności otrzymanej funkcji oraz wyznacza równania asymptot jej wykresu; \\
		
		\textbullet\ sprawdza, czy dany punkt należy do wykresu funkcji $f(x)=\frac{a}{x}$; \\
		
		\textbullet\ wyznacza dziedzinę prostego wyrażenia wymiernego; \\
		
		\textbullet\ oblicza w prostych przypadkach wartość wyrażenia wymiernego dla danej wartości zmiennej; \\
		
		\textbullet\ upraszcza wyrażenia wymierne w przypadkach wymagających wyłączenia czynnika liczbowego przed nawias w liczniku lub mianowniku; \\
		
		\textbullet\ mnoży i dzieli wyrażenia wymierne w prostych przypadkach oraz podaje odpowiednie założenia; \\
		
		\textbullet\ rozwiązuje równania wymierne postaci $\frac{u(x)}{w(x)}=0$, podając i uwzględniając odpowiednie założenia; \\
		
		\textbullet\ stosuje własności wartości bezwzględnej do rozwiązywania równań postaci $|x|=b$. \\ \hline
		
		\textbf{Wymagania edukacyjne na ocenę dostateczną} \\ \hline
		
		\textbullet\ podaje własności funkcji $f(x)=\frac{a}{x}+q$, gdzie $a\neq 0$, oraz wyznacza równania asymptot jej wykresu; \\
		
		\textbullet\ podaje własności funkcji $f(x)=\frac{a}{x-p}$, gdzie $a\neq 0$, oraz wyznacza równania asymptot jej wykresu; \\
		
		\textbullet\ sprawdza, czy dany punkt należy do wykresu funkcji $f(x)=\frac{a}{x-p}+q$; \\
		
		\textbullet\ wyznacza wartość parametru $q$ we wzorze funkcji, np. $f(x)=\frac{5}{x}+q$, gdy dane są współrzędne punktu należącego do jej wykresu; \\
		
		\textbullet\ wyznacza dziedzinę wyrażenia wymiernego; \\
		
		\textbullet\ oblicza wartość wyrażenia wymiernego dla danej wartości zmiennej; \\
		
		\textbullet\ upraszcza wyrażenia wymierne, stosując wzory skróconego mnożenia do rozkładu licznika lub mianownika na czynniki; \\
		
		\textbullet\ mnoży i dzieli wyrażenia wymierne oraz podaje odpowiednie założenia; \\
		
		\textbullet\ rozwiązuje równania wymierne postaci $\frac{u(x)}{w(x)}=a$, gdzie $a\neq 0$, podając i uwzględniając odpowiednie założenia; \\
		
		\textbullet\ wykorzystuje wyrażenia wymierne do rozwiązywania prostych zadań tekstowych; \\
		
		\textbullet\ stosuje własności wartości bezwzględnej do rozwiązywania równań postaci $|x-a|=b$. \\ \hline
		
		\textbf{Wymagania edukacyjne na ocenę dobrą} \\ \hline
		
		\textbullet\ szkicuje wykres funkcji $f(x)=\frac{a}{x}$, gdzie $a\neq 0$, w podanym zbiorze oraz podaje jej własności: dziedzinę, zbiór wartości i przedziały monotoniczności; \\
		
		\textbullet\ dobiera wzór funkcji do jej wykresu; \\
		
		\textbullet\ wyznacza współczynnik $a$ tak, aby funkcja $f(x)=\frac{a}{x}$ spełniała podane warunki; \\
		
		\textbullet\ szkicuje wykres funkcji $f(x)=\frac{a}{x-p}+q$, gdzie $x\in\mathbb{R}\setminus\{p\}$ i $a\neq 0$, oraz wyznacza równania jego asymptot; \\
		
		\textbullet\ mnoży i dzieli wyrażenia wymierne w trudniejszych przypadkach oraz podaje odpowiednie założenia; \\
		
		\textbullet\ określa dziedzinę funkcji, której wzór zawiera ułamek; \\
		
		\textbullet\ przekształca wzór, wyznaczając z niego wskazaną zmienną; \\
		
		\textbullet\ rozwiązuje równania wymierne w trudniejszych przypadkach; \\
		
		\textbullet\ podaje interpretację geometryczną rozwiązania równania wymiernego; \\
		
		\textbullet\ wykorzystuje wyrażenia wymierne do rozwiązywania zadań tekstowych. \\ \hline
		
		\textbf{Wymagania edukacyjne na ocenę bardzo dobrą} \\ \hline
		
		\textbullet\ szkicuje wykres funkcji $f(x)=\frac{a}{x}$, gdzie $a\neq 0$, w podanym zbiorze w trudniejszych przypadkach; \\
		
		\textbullet\ wyznacza wzór funkcji spełniającej podane warunki; \\
		
		\textbullet\ wykorzystuje wyrażenia wymierne do rozwiązywania trudniejszych zadań tekstowych; \\
		
		\textbullet\ stosuje własności wartości bezwzględnej do rozwiązywania równań w trudniejszych przypadkach. \\ \hline
		
		\textbf{Wymagania edukacyjne na ocenę celującą} \\ \hline
		
		\textbullet\ stosuje funkcje i wyrażenia wymierne do rozwiązywania zadań o podwyższonym stopniu trudności. \\ \hline
		
	\end{longtable}
	
	\section{Planimetria}
	
	\begin{longtable}{|p{16cm}|}
		\hline
		
		\textbf{Wymagania edukacyjne na ocenę dopuszczającą} \\ \hline
		
		\textbullet\ rozróżnia trójkąty ostrokątne, prostokątne i rozwartokątne; \\
		
		\textbullet\ stosuje twierdzenie o sumie miar kątów w trójkącie w prostych przypadkach; \\
		
		\textbullet\ uzasadnia przystawanie trójkątów, wykorzystując cechy przystawania, w prostych przypadkach; \\
		
		\textbullet\ zapisuje proporcje długości odpowiednich boków trójkątów podobnych; \\
		
		\textbullet\ sprawdza, czy dane figury są podobne; \\
		
		\textbullet\ wskazuje w wielokątach odcinki, których długości są proporcjonalne. \\ \hline
		
		\textbf{Wymagania edukacyjne na ocenę dostateczną} \\ \hline
		
		\textbullet\ sprawdza, czy z trzech odcinków o danych długościach można zbudować trójkąt; \\
		
		\textbullet\ uzasadnia podobieństwo trójkątów, wykorzystując cechy podobieństwa, w prostych przypadkach; \\
		
		\textbullet\ wykorzystuje podobieństwo trójkątów do rozwiązywania prostych zadań; \\
		
		\textbullet\ stosuje podobieństwo wielokątów do obliczania długości boków; \\
		
		\textbullet\ stosuje twierdzenie o stosunku pól figur podobnych w prostych zadaniach; \\
		
		\textbullet\ rozwiązuje proste zadania, wykorzystując twierdzenie Talesa. \\ \hline
		
		\textbf{Wymagania edukacyjne na ocenę dobrą} \\ \hline
		
		\textbullet\ przeprowadza dowód twierdzenia o sumie miar kątów w trójkącie; \\
		
		\textbullet\ stosuje cechy przystawania trójkątów do rozwiązywania zadań geometrycznych; \\
		
		\textbullet\ wykorzystuje podobieństwo trójkątów do rozwiązywania problemów praktycznych i zadań geometrycznych; \\
		
		\textbullet\ rozwiązuje zadania dotyczące podobieństwa wielokątów; \\
		
		\textbullet\ rozwiązuje zadania, wykorzystując twierdzenie Talesa; \\
		
		\textbullet\ rozwiązuje zadania wymagające uzasadnienia lub dowodu, korzystając z przystawania trójkątów. \\ \hline
		
		\textbf{Wymagania edukacyjne na ocenę bardzo dobrą} \\ \hline
		
		\textbullet\ przeprowadza dowód twierdzenia o mierze kąta zewnętrznego trójkąta; \\
		
		\textbullet\ udowadnia, że symetralne boków trójkąta przecinają się w jednym punkcie; \\
		
		\textbullet\ udowadnia, że dwusieczne kątów trójkąta przecinają się w jednym punkcie; \\
		
		\textbullet\ stosuje cechy przystawania trójkątów do rozwiązywania trudniejszych zadań geometrycznych; \\
		
		\textbullet\ wykorzystuje podobieństwo trójkątów do rozwiązywania problemów praktycznych i trudniejszych zadań geometrycznych; \\
		
		\textbullet\ rozwiązuje zadania wymagające uzasadnienia lub dowodu, korzystając z podobieństwa trójkątów; \\
		
		\textbullet\ rozwiązuje zadania wymagające uzasadnienia lub dowodu z zastosowaniem twierdzenia Talesa. \\ \hline
		
		\textbf{Wymagania edukacyjne na ocenę celującą} \\ \hline
		
		\textbullet\ przeprowadza dowód twierdzenia Talesa; \\
		
		\textbullet\ rozwiązuje zadania o znacznym stopniu trudności dotyczące przystawania i podobieństwa figur. \\ \hline
		
	\end{longtable}
	
	\section{Trygonometria}
	
	\begin{longtable}{|p{16cm}|}
		\hline
		
		\textbf{Wymagania edukacyjne na ocenę dopuszczającą} \\ \hline
		
		\textbullet\ stosuje w prostych przypadkach twierdzenie Pitagorasa i twierdzenie odwrotne do twierdzenia Pitagorasa; \\
		
		\textbullet\ wykorzystuje wzory na długość przekątnej kwadratu i wysokość trójkąta równobocznego; \\
		
		\textbullet\ oblicza wartości funkcji trygonometrycznych kąta ostrego w trójkącie prostokątnym o danych długościach boków; \\
		
		\textbullet\ odczytuje z tablic wartości funkcji trygonometrycznych danego kąta ostrego; \\
		
		\textbullet\ odczytuje z tablic miarę kąta ostrego, gdy dana jest wartość jego funkcji trygonometrycznej; \\
		
		\textbullet\ rozwiązuje trójkąty prostokątne w prostych przypadkach; \\
		
		\textbullet\ przedstawia kąt wypukły na rysunku i oblicza wartości jego funkcji trygonometrycznych, gdy dane są współrzędne punktu leżącego na jego końcowym ramieniu; \\
		
		\textbullet\ wyznacza, korzystając z definicji, wartości funkcji trygonometrycznych kątów $0^\circ$, $90^\circ$ i $180^\circ$; \\
		
		\textbullet\ stosuje w zadaniach wzór na pole trójkąta $P=\frac{1}{2}ah$ oraz wzór na pole trójkąta równobocznego o boku $a$: $P=\frac{a^2\sqrt{3}}{4}$; \\
		
		\textbullet\ rozróżnia czworokąty: kwadrat, prostokąt, romb, równoległobok i trapez oraz zna ich własności. \\ \hline
		
		\textbf{Wymagania edukacyjne na ocenę dostateczną} \\ \hline
		
		\textbullet\ podaje wartości funkcji trygonometrycznych kątów $30^\circ$, $45^\circ$ i $60^\circ$; \\
		
		\textbullet\ podaje związki między funkcjami trygonometrycznymi tego samego kąta; \\
		
		\textbullet\ oblicza wartości pozostałych funkcji trygonometrycznych kąta, gdy dana jest wartość jego sinusa lub cosinusa; \\
		
		\textbullet\ rozwiązuje trójkąty prostokątne; \\
		
		\textbullet\ stosuje funkcje trygonometryczne do rozwiązywania prostych zadań osadzonych w kontekście praktycznym; \\
		
		\textbullet\ oblicza wartości funkcji trygonometrycznych kątów $120^\circ$, $135^\circ$ i $150^\circ$; \\
		
		\textbullet\ stosuje wzory $\sin(180^\circ-\alpha)=\sin\alpha$, $\cos(180^\circ-\alpha)=-\cos\alpha$ i $\tan(180^\circ-\alpha)=-\tan\alpha$ do obliczania wartości funkcji trygonometrycznych; \\
		
		\textbullet\ oblicza wartości funkcji trygonometrycznych kątów rozwartych, korzystając z tablic wartości funkcji trygonometrycznych; \\
		
		\textbullet\ wykorzystuje funkcje trygonometryczne do obliczania obwodów i pól podstawowych czworokątów w prostych przypadkach. \\ \hline
		
		\textbf{Wymagania edukacyjne na ocenę dobrą} \\ \hline
		
		\textbullet\ wyznacza długości odcinków w trójkącie, korzystając z twierdzenia Pitagorasa; \\
		
		\textbullet\ wyprowadza ogólne zależności, np. wzory na długość przekątnej kwadratu i wysokość trójkąta równobocznego; \\
		
		\textbullet\ wyznacza wartości funkcji trygonometrycznych kątów ostrych w bardziej złożonych sytuacjach; \\
		
		\textbullet\ uzasadnia proste zależności, korzystając z własności funkcji trygonometrycznych; \\
		
		\textbullet\ stosuje funkcje trygonometryczne do rozwiązywania trójkątów w zadaniach osadzonych w kontekście praktycznym; \\
		
		\textbullet\ uzasadnia związki między funkcjami trygonometrycznymi kątów ostrych $\alpha$ i $90^\circ-\alpha$; \\
		
		\textbullet\ przekształca w prostych przypadkach wyrażenia trygonometryczne, wykorzystując związki między funkcjami trygonometrycznymi tego samego kąta; \\
		
		\textbullet\ oblicza wartości pozostałych funkcji trygonometrycznych kąta, gdy dana jest wartość jego tangensa; \\
		
		\textbullet\ rysuje kąt wypukły w układzie współrzędnych, gdy dana jest wartość jego tangensa; \\
		
		\textbullet\ stosuje w zadaniach wzory na pole trójkąta, w tym wzór $P=\frac{1}{2}ab\sin\gamma$; \\
		
		\textbullet\ oblicza pola czworokątów; \\
		
		\textbullet\ wykorzystuje umiejętność obliczania pól trójkątów do wyznaczania pól innych wielokątów. \\ \hline
		
		\textbf{Wymagania edukacyjne na ocenę bardzo dobrą} \\ \hline
		
		\textbullet\ stosuje poznane związki między funkcjami trygonometrycznymi do upraszczania wyrażeń trygonometrycznych; \\
		
		\textbullet\ wyprowadza wzór na jedynkę trygonometryczną oraz pozostałe związki między funkcjami trygonometrycznymi tego samego kąta; \\
		
		\textbullet\ przekształca wyrażenia trygonometryczne, wykorzystując związki między funkcjami trygonometrycznymi tego samego kąta; \\
		
		\textbullet\ stosuje wzór na pole trójkąta $P=\frac{1}{2}ab\sin\gamma$ w zadaniach o podwyższonym stopniu trudności; \\
		
		\textbullet\ stosuje wzór Herona do obliczania pola trójkąta; \\
		
		\textbullet\ oblicza pola czworokątów w trudniejszych przypadkach; \\
		
		\textbullet\ wykorzystuje umiejętność obliczania pól trójkątów do wyznaczania pól innych wielokątów; \\
		
		\textbullet\ uzasadnia zależności miarowe w czworokątach; \\
		
		\textbullet\ dowodzi wzoru $P=\frac{1}{2}ab\sin\gamma$. \\ \hline
		
		\textbf{Wymagania edukacyjne na ocenę celującą} \\ \hline
		
		\textbullet\ przeprowadza dowody twierdzenia Pitagorasa i twierdzenia odwrotnego do twierdzenia Pitagorasa; \\
		
		\textbullet\ rozwiązuje zadania o znacznym stopniu trudności z zastosowaniem trygonometrii, w tym zadania wymagające dowodzenia zależności miarowych w trójkątach i czworokątach. \\ \hline
		
	\end{longtable}
	
	\section{Planimetria}
	
	\begin{longtable}{|p{16cm}|}
		\hline
		
		\textbf{Wymagania edukacyjne na ocenę dopuszczającą} \\ \hline
		
		\textbullet\ rozpoznaje kąty środkowe w okręgu; \\
		
		\textbullet\ oblicza długość okręgu i długość łuku okręgu w prostych przypadkach; \\
		
		\textbullet\ oblicza pole koła i pole wycinka koła w prostych przypadkach; \\
		
		\textbullet\ rozpoznaje okręgi styczne, gdy dane są ich promienie i odległość między ich środkami; \\
		
		\textbullet\ rozpoznaje kąty wpisane w okrąg oraz wskazuje łuki, na których są oparte; \\
		
		\textbullet\ stosuje w prostych przypadkach twierdzenie o kącie środkowym i kącie wpisanym opartych na tym samym łuku oraz wnioski z tego twierdzenia; \\
		
		\textbullet\ rozwiązuje zadania dotyczące okręgu opisanego na trójkącie równobocznym lub prostokątnym; \\
		
		\textbullet\ rozwiązuje zadania dotyczące okręgu wpisanego w trójkąt równoboczny lub prostokątny; \\
		
		\textbullet\ opisuje własności wielokątów foremnych; \\
		
		\textbullet\ stosuje w prostych przypadkach twierdzenie cosinusów do rozwiązywania trójkątów; \\
		
		\textbullet\ wskazuje najmniejszy lub największy kąt w trójkącie, gdy dane są długości jego boków. \\ \hline
		
		\textbf{Wymagania edukacyjne na ocenę dostateczną} \\ \hline
		
		\textbullet\ oblicza pole figury, stosując wzory na pole koła i pole wycinka koła w prostych sytuacjach; \\
		
		\textbullet\ wykorzystuje twierdzenie o kącie środkowym i kącie wpisanym opartych na tym samym łuku oraz wnioski z tego twierdzenia; \\
		
		\textbullet\ rozwiązuje proste zadania dotyczące okręgu opisanego na dowolnym trójkącie; \\
		
		\textbullet\ rozwiązuje proste zadania dotyczące okręgu wpisanego w dowolny trójkąt; \\
		
		\textbullet\ oblicza miarę kąta wewnętrznego wielokąta foremnego; \\
		
		\textbullet\ wyznacza liczbę boków wielokąta foremnego, gdy dana jest suma miar jego kątów wewnętrznych; \\
		
		\textbullet\ oblicza w prostych przypadkach promień okręgu opisanego na wielokącie foremnym i promień okręgu wpisanego w wielokąt foremny; \\
		
		\textbullet\ stosuje twierdzenie cosinusów do rozwiązywania trójkątów. \\ \hline
		
		\textbf{Wymagania edukacyjne na ocenę dobrą} \\ \hline
		
		\textbullet\ oblicza pole figury, stosując wzory na pole koła i pole wycinka koła; \\
		
		\textbullet\ wykorzystuje twierdzenie o odcinkach stycznych do rozwiązywania zadań; \\
		
		\textbullet\ stosuje w trudniejszych przypadkach twierdzenie o kącie środkowym i kącie wpisanym opartych na tym samym łuku oraz wnioski z tego twierdzenia; \\
		
		\textbullet\ stosuje twierdzenie o kącie między styczną a cięciwą okręgu do rozwiązywania zadań; \\
		
		\textbullet\ stosuje wzory $P=\frac{abc}{4R}$ i $P=\frac{a+b+c}{2}\cdot r$ do obliczania pola trójkąta; \\
		
		\textbullet\ określa, czy trójkąt jest ostrokątny, prostokątny czy rozwartokątny, gdy dane są długości jego boków; \\
		
		\textbullet\ rozwiązuje zadania dotyczące okręgu opisanego na trójkącie; \\
		
		\textbullet\ rozwiązuje zadania dotyczące okręgu wpisanego w trójkąt; \\
		
		\textbullet\ stosuje twierdzenie cosinusów do rozwiązywania trójkątów oraz zadań osadzonych w kontekście praktycznym. \\ \hline
		
		\textbf{Wymagania edukacyjne na ocenę bardzo dobrą} \\ \hline
		
		\textbullet\ wykorzystuje twierdzenie o odcinkach stycznych do rozwiązywania trudniejszych zadań; \\
		
		\textbullet\ stosuje twierdzenie o kącie między styczną a cięciwą okręgu do rozwiązywania trudniejszych zadań; \\
		
		\textbullet\ uzasadnia wzory $P=\frac{abc}{4R}$ i $P=\frac{a+b+c}{2}\cdot r$ na pole trójkąta; \\
		
		\textbullet\ stosuje twierdzenie cosinusów do rozwiązywania trójkątów i zadań osadzonych w kontekście praktycznym; \\
		
		\textbullet\ udowadnia, że symetralne boków trójkąta przecinają się w jednym punkcie; \\
		
		\textbullet\ udowadnia, że dwusieczne kątów wewnętrznych trójkąta przecinają się w jednym punkcie. \\ \hline
		
		\textbf{Wymagania edukacyjne na ocenę celującą} \\ \hline
		
		\textbullet\ udowadnia zależności dotyczące wielokątów foremnych w zadaniach o podwyższonym stopniu trudności; \\
		
		\textbullet\ konstruuje pięciokąt foremny; \\
		
		\textbullet\ przeprowadza dowody twierdzeń o kącie środkowym i kącie wpisanym oraz o kątach wpisanych opartych na tym samym łuku; \\
		
		\textbullet\ przeprowadza dowód twierdzenia o cięciwach w okręgu; \\
		
		\textbullet\ przeprowadza dowód twierdzenia cosinusów; \\
		
		\textbullet\ rozwiązuje zadania o podwyższonym stopniu trudności z zastosowaniem trygonometrii. \\ \hline
		
	\end{longtable}
	
\end{document}